\section{The Wallfacer Experiment}

\subsection{Motivation and Setup}

The experiment was motivated by a straightforward question: what happens when an LLM-based agent pipeline is left to operate autonomously for an extended period? While single-turn and short-session evaluations of LLM agents are common in the literature, sustained autonomous operation over days or weeks remains underexplored. We sought to observe, rather than prescribe, the architectural pressures that emerge at scale.

The Wallfacer pipeline was deployed on a virtual private server running Ubuntu, with Claude Code \cite{anthropic2024claude} as the base LLM substrate. The pipeline was configured for a software engineering task: maintaining and extending a moderately complex codebase. It operated in a self-directed loop encompassing ideation (selecting the next task), implementation (writing code), testing (running the test suite), documentation (updating project records), and deployment (committing and pushing changes). No human reviewed or intervened in any of these stages for the duration of the experiment.

We conducted the experiment across four successive architectural configurations, each introduced when the preceding configuration exhibited sustained performance degradation that we judged, in retrospect, to stem from a structural rather than parametric limitation. The four configurations were: (1) a single agent performing all functions in a self-loop; (2) a Strategist agent separated from an Executor agent; (3) the addition of a Cartographer agent; and (4) the further addition of a Critic agent.\footnote{We name these roles after their cognitive mode rather than their mechanical function. The Strategist decides \emph{what} to do (not merely ``proposes''), the Executor implements (not merely ``builds''), the Cartographer produces legible abstractions of system state (not merely ``documents''), and the Critic exercises evaluative judgment (not merely ``verifies''). These names reflect the fission principle's emphasis on competing cognitive modes.} Each configuration was run until we observed a clear and persistent bottleneck, at which point we introduced the next configuration. The total duration across all configurations was seven days, during which the codebase grew from approximately 10,000 to over 40,000 lines of code.

\subsection{Observations and Analysis}

\begin{figure*}[t]
\centering
\begin{tikzpicture}[
    agent/.style={rectangle, rounded corners=3pt, draw=#1!70, fill=#1!12, minimum width=1.3cm, minimum height=0.55cm, font=\scriptsize\sffamily, inner sep=2pt, align=center},
    agent/.default={cBuilder},
    phase/.style={font=\small\sffamily\bfseries, text=black!80},
    bottleneck/.style={font=\tiny\sffamily\itshape, text=cBottleneck, align=center},
    arr/.style={-{Stealth[length=4pt]}, thick, black!50},
    farr/.style={-{Stealth[length=3pt]}, thick, dashed, black!35},
    trans/.style={-{Stealth[length=5pt]}, very thick, black!20},
]

\def\sep{3.8}

% Phase 1
\node[phase] at (0, 2.5) {Phase 1};
\node[agent=cBuilder, minimum width=2cm] (b1) at (0, 1.5) {Executor\\[-1pt]{\tiny(all functions)}};
\draw[arr, looseness=5] (b1.160) to[out=110,in=70] (b1.20);
\node[bottleneck] at (0, 0.3) {context\\[-1pt]competition};

\draw[trans] (1.2, 1.0) -- (\sep-1.2, 1.0);

% Phase 2
\node[phase] at (\sep, 2.5) {Phase 2};
\node[agent=cProposer] (p2) at (\sep, 2.0) {Strategist};
\node[agent=cBuilder] (b2) at (\sep, 1.0) {Executor};
\draw[arr] (p2.south) -- (b2.north);
\draw[farr] (b2.west) to[out=170,in=190] (p2.west);
\node[bottleneck] at (\sep, 0.3) {knowledge\\[-1pt]decay};

\draw[trans] (\sep+1.4, 1.5) -- (2*\sep-1.4, 1.5);

% Phase 3
\node[phase] at (2*\sep, 2.5) {Phase 3};
\node[agent=cProposer] (p3) at (2*\sep-0.9, 2.0) {Strategist};
\node[agent=cBuilder] (b3) at (2*\sep+0.9, 2.0) {Executor};
\node[agent=cDocumenter] (d3) at (2*\sep, 0.85) {Cartographer};
\draw[arr] (p3) -- (b3);
\draw[arr] (b3.south) -- (d3.north east);
\draw[farr] (d3.north west) -- (p3.south);
\node[bottleneck] at (2*\sep, 0.1) {quality\\[-1pt]control};

\draw[trans] (2*\sep+1.4, 1.5) -- (3*\sep-1.4, 1.5);

% Phase 4
\node[phase] at (3*\sep, 2.5) {Phase 4};
\node[agent=cProposer] (p4) at (3*\sep-0.9, 2.0) {Strategist};
\node[agent=cBuilder] (b4) at (3*\sep+0.9, 2.0) {Executor};
\node[agent=cDocumenter] (d4) at (3*\sep-0.9, 0.85) {Cart.};
\node[agent=cVerifier] (v4) at (3*\sep+0.9, 0.85) {Critic};
\draw[arr] (p4) -- (b4);
\draw[arr] (b4.south) -- (v4.north);
\draw[arr] (v4) -- (d4);
\draw[farr] (d4.north) -- (p4.south);
\draw[farr, cVerifier!60] (v4.east) to[out=0,in=0,looseness=1.8] (b4.east);
\node[bottleneck] at (3*\sep, 0.1) {coordination\\[-1pt]overhead};

% Bottom annotation
\draw[decorate, decoration={brace, amplitude=5pt, mirror}, thick, black!30] (-1.3, -0.3) -- (3*\sep+1.3, -0.3);
\node[font=\tiny\sffamily\itshape, text=black!50] at (1.5*\sep, -0.8) {Empirically observed in the Wallfacer experiment};

\end{tikzpicture}
\caption{Architectural evolution across four observed phases. Each transition was triggered by a persistent bottleneck (shown in red). Solid arrows indicate primary information flow; dashed arrows indicate feedback.}
\label{fig:phases}
\end{figure*}

We report three principal observations from the experiment, along with our interpretation of each. We acknowledge that as a single case study, these observations are suggestive rather than conclusive, and that alternative interpretations are possible.

\textbf{Observation 1: Bottleneck migration.} Across all four configurations, resolving a performance bottleneck in one dimension of system behavior reliably relocated the binding constraint to a different dimension, rather than producing an overall improvement that persisted. For example, in the two-agent configuration, separating planning from implementation eliminated the context-competition bottleneck observed in the single-agent configuration, but introduced a new bottleneck: implicit knowledge generated during implementation (such as the rationale behind design decisions) was not captured in any durable form and therefore decayed between planning cycles. The Cartographer was introduced to address this decay, which it did, but at the cost of introducing a quality-control bottleneck: three agents were producing outputs, but none was tasked with evaluating whether those outputs met any standard of correctness.

This pattern, in which resolving one constraint exposes another, aligns with Goldratt's Theory of Constraints \cite{goldratt1984}, which posits that any system has exactly one binding constraint at a time and that improving non-constraints yields no net throughput gain. It also resonates with the No Free Lunch theorems \cite{wolpert1997}, which establish that no single optimization strategy dominates across all problem instances. In our setting, the analogous claim is that no single architectural configuration dominates across all bottleneck types. We note, however, that our evidence for this claim is limited to a single system operating on a single class of tasks, and that broader empirical validation is required before any general principle can be asserted with confidence.

\textbf{Observation 2: Emergent self-diagnosis.} In the three-agent

\begin{figure}[t]
\centering
\begin{tikzpicture}[scale=0.82,
    cell/.style={minimum width=1.55cm, minimum height=0.34cm, font=\tiny\sffamily, inner sep=1pt},
    hi/.style={cell, fill=cVerifier!20, draw=cVerifier!50},
    lo/.style={cell, fill=cEmpirical!12, draw=cEmpirical!30},
    md/.style={cell, fill=cConjectured!10, draw=cConjectured!30},
    lbl/.style={font=\tiny\sffamily\bfseries, text=black!70, anchor=east},
    plbl/.style={font=\scriptsize\sffamily\bfseries, text=black!80}
]
\def\rh{0.48}
\def\cw{1.75}

\node[plbl] at (0, 0) {Ph.1};
\node[plbl] at (\cw, 0) {Ph.2};
\node[plbl] at (2*\cw, 0) {Ph.3};

\node[lbl] at (-0.9, -\rh) {Context};
\node[lbl] at (-0.9, -2*\rh) {Knowl.};
\node[lbl] at (-0.9, -3*\rh) {Quality};
\node[lbl] at (-0.9, -4*\rh) {Coord.};

\node[hi] at (0, -\rh) {\textbf{binding}};
\node[lo] at (0, -2*\rh) {};
\node[lo] at (0, -3*\rh) {};
\node[lo] at (0, -4*\rh) {};

\node[lo] at (\cw, -\rh) {};
\node[hi] at (\cw, -2*\rh) {\textbf{binding}};
\node[lo] at (\cw, -3*\rh) {};
\node[md] at (\cw, -4*\rh) {};

\node[lo] at (2*\cw, -\rh) {};
\node[lo] at (2*\cw, -2*\rh) {};
\node[hi] at (2*\cw, -3*\rh) {\textbf{binding}};
\node[md] at (2*\cw, -4*\rh) {};

\draw[-{Stealth[length=3pt]}, thick, cVerifier!60, dashed]
    (0.5, -\rh) to[out=-15, in=165] (\cw-0.5, -2*\rh);
\draw[-{Stealth[length=3pt]}, thick, cVerifier!60, dashed]
    (\cw+0.5, -2*\rh) to[out=-15, in=165] (2*\cw-0.5, -3*\rh);

\end{tikzpicture}
\caption{Bottleneck migration across phases. Dark cells mark the binding constraint; light cells indicate resolved dimensions. Dashed arrows show how each fission relocates the bottleneck.}
\label{fig:migration}
\end{figure}

\begin{figure*}[t]
\centering
\begin{tikzpicture}[x=1.5cm, y=3.2cm]
% Axes
\draw[-{Stealth[length=4pt]}, black!50] (0,0) -- (8.5,0) node[right, font=\tiny\sffamily, xshift=2pt] {Phase};
\draw[-{Stealth[length=4pt]}, black!50] (0,0) -- (0,1.18) node[above, font=\tiny\sffamily] {\%};

% Y-axis ticks
\foreach \y/\l in {0/0, 0.25/25, 0.5/50, 0.75/75, 1.0/100} {
    \draw[black!20] (0,\y) -- (8.2,\y);
    \node[left, font=\tiny\sffamily, text=black!50] at (-0.1,\y) {\l};
}

% X-axis labels
\foreach \x/\l in {1/1, 2/2, 3/3, 4/4, 5/5, 6/6, 7/7, 8/8} {
    \node[below, font=\tiny\sffamily, text=black!60] at (\x,-0.04) {\l};
}

% Empirical / conjectured boundary - subtle, no text in data area
\draw[dashed, black!20, thin] (4.5,-0.05) -- (4.5,1.02);
\node[font=\tiny\sffamily\itshape, text=black!30, anchor=north] at (2.25,-0.12) {empirical};
\node[font=\tiny\sffamily\itshape, text=black!30, anchor=north] at (6.35,-0.12) {conjectured};

% Data: Capability (teal)
\draw[cEmpirical, thick] plot coordinates {(1,0.60) (2,0.75) (3,0.85) (4,0.92) (5,0.95) (6,0.96) (7,0.98) (8,0.99)};
\foreach \x/\y in {1/0.60, 2/0.75, 3/0.85, 4/0.92, 5/0.95, 6/0.96, 7/0.98, 8/0.99} {
    \fill[cEmpirical] (\x,\y) circle (1.5pt);
}

% Data: Coordination cost (red)
\draw[cVerifier, thick] plot coordinates {(1,0.05) (2,0.15) (3,0.25) (4,0.38) (5,0.48) (6,0.52) (7,0.62) (8,0.58)};
\foreach \x/\y in {1/0.05, 2/0.15, 3/0.25, 4/0.38, 5/0.48, 6/0.52, 7/0.62, 8/0.58} {
    \fill[cVerifier] (\x,\y) circle (1.5pt);
}

% Data: Self-awareness (purple)
\draw[cCoordinator, thick] plot coordinates {(1,0.00) (2,0.00) (3,0.05) (4,0.08) (5,0.15) (6,0.30) (7,0.40) (8,0.80)};
\foreach \x/\y in {1/0.00, 2/0.00, 3/0.05, 4/0.08, 5/0.15, 6/0.30, 7/0.40, 8/0.80} {
    \fill[cCoordinator] (\x,\y) circle (1.5pt);
}

% Data: Boundary fluidity (orange)
\draw[cDocumenter, thick, densely dashed] plot coordinates {(1,0.00) (2,0.00) (3,0.00) (4,0.00) (5,0.02) (6,0.03) (7,0.04) (8,0.05)};
\foreach \x/\y in {5/0.02, 6/0.03, 7/0.04, 8/0.05} {
    \fill[cDocumenter] (\x,\y) circle (1.2pt);
}

% Data: Autonomy (blue)
\draw[cBuilder, thick, densely dashed] plot coordinates {(1,0.10) (2,0.15) (3,0.20) (4,0.25) (5,0.30) (6,0.35) (7,0.40) (8,0.50)};
\foreach \x/\y in {1/0.10, 2/0.15, 3/0.20, 4/0.25, 5/0.30, 6/0.35, 7/0.40, 8/0.50} {
    \fill[cBuilder] (\x,\y) circle (1.2pt);
}

% Data: Ontological depth (pink/monitor green)
\draw[cMonitor, thick, densely dotted] plot coordinates {(1,0.00) (2,0.00) (3,0.02) (4,0.05) (5,0.08) (6,0.12) (7,0.18) (8,0.25)};
\foreach \x/\y in {5/0.08, 6/0.12, 7/0.18, 8/0.25} {
    \fill[cMonitor] (\x,\y) circle (1.2pt);
}

% Legend
\node[font=\tiny\sffamily, anchor=west] at (0.2,1.08) {%
    \textcolor{cEmpirical}{\rule{6pt}{2pt}} Capability\quad
    \textcolor{cVerifier}{\rule{6pt}{2pt}} Coordination cost\quad
    \textcolor{cCoordinator}{\rule{6pt}{2pt}} Self-awareness\quad
    \textcolor{cBuilder}{\rule[0.3ex]{6pt}{0.4pt}\rule[0.3ex]{2pt}{0.4pt}\rule[0.3ex]{6pt}{0.4pt}} Autonomy\quad
    \textcolor{cDocumenter}{\rule[0.3ex]{6pt}{0.4pt}\rule[0.3ex]{2pt}{0.4pt}\rule[0.3ex]{6pt}{0.4pt}} Boundary fluidity\quad
    \textcolor{cMonitor}{$\cdots$} Ontological depth
};

% Annotations - right margin, outside plot area
\draw[-{Stealth[length=3pt]}, cCoordinator!50, thin] (8.6,0.88) -- (8.05,0.81);
\node[font=\tiny\sffamily, text=cCoordinator!80, anchor=west, align=left] at (8.65,0.88) {self-awareness\\[-1pt]jump};

\draw[-{Stealth[length=3pt]}, cVerifier!50, thin] (8.6,0.52) -- (8.05,0.58);
\node[font=\tiny\sffamily, text=cVerifier!80, anchor=west, align=left] at (8.65,0.52) {cost\\[-1pt]declines};

\end{tikzpicture}
\caption{System metric evolution across eight phases. Solid lines represent primary metrics (capability, coordination cost, self-awareness); dashed and dotted lines represent emergent metrics that become significant only in later phases (autonomy, boundary fluidity, ontological depth). Phases 1 to 4 are empirically grounded; Phases 5 to 8 are conjectured extrapolations. The sharp rise in self-awareness at Phase 8 reflects the Architect role's capacity for architectural self-modification. Values are normalized estimates based on the Wallfacer observations (Phases 1--4) and theoretical projections.}
\label{fig:metrics}
\end{figure*}

configuration (Strategist, Executor, Cartographer), the system proposed refactoring its own monolithic codebase at approximately 51,000 lines of code. This behavior did not emerge in the two-agent configuration at a comparable scale. We interpret this as evidence that the Cartographer role, by producing structured representations of the system's architecture and history, provided the Strategist with a legible abstraction of global system state that was otherwise unavailable. In the two-agent configuration, the Strategist's decisions were conditioned only on the current codebase and task queue, without any higher-level summary of architectural trends.

This interpretation finds support in Clark and Chalmers' Extended Mind hypothesis \cite{clark1998}, which argues that cognitive processes are not confined to the brain but extend to external artifacts that participate in reasoning. In our case, the documentation artifacts produced by the Cartographer function as external cognitive scaffolding for the Strategist: they do not merely record what has happened but actively shape what the Strategist can perceive and reason about. We are cautious, however, about overinterpreting a single instance of emergent behavior. It is possible that the refactoring proposal would have emerged eventually in the two-agent configuration at a higher line count, or that it was triggered by idiosyncratic features of the codebase rather than by the architectural difference we hypothesize.

\textbf{Observation 3: Coordination cost scaling.} As the number of agent roles increased from one to four, we observed a nonlinear increase in the volume of inter-agent communication required to maintain coherent system behavior. The four-agent configuration exhibited episodes of circular dependency, in which the Critic rejected output that the Executor had produced according to the Strategist's specification, leading to repeated revision cycles that consumed computational resources without advancing the project. This observation is consistent with Brooks' \cite{brooks1975} well-known finding that communication overhead in human teams scales quadratically with team size, and it suggests that the benefits of role differentiation are subject to diminishing returns as coordination costs accumulate. Figure~\ref{fig:metrics} traces the estimated trajectory of six system metrics across all eight phases, illustrating how capability saturates while coordination cost, self-awareness, and autonomy follow distinct growth curves.
